\documentclass[11pt]{amsart}

\usepackage[T1]{fontenc}
\usepackage{lmodern}
\usepackage{microtype}
\usepackage{amsmath,amssymb,amsthm}
\usepackage{booktabs}
\usepackage[hidelinks]{hyperref}


\hypersetup{
  pdftitle={The Ivan--Wang Conjecture on Induced N-Saturation for 1 <= n <= 5}
}

\newtheorem{theorem}{Theorem}
\newtheorem{lemma}[theorem]{Lemma}
\newtheorem{proposition}[theorem]{Proposition}

\newcommand{\Bn}{\mathcal B_n}
\newcommand{\Nposet}{\mathcal N}
\newcommand{\satstar}{\operatorname{sat}^{*}}

\title[Induced $\mathcal N$-saturation for $n\le5$]{The Ivan--Wang Conjecture on Induced
\texorpdfstring{$\mathcal N$}{N}-Saturation for
\texorpdfstring{$1\leq n\leq 5$}{1 <= n <= 5}}
\date{}
\subjclass[2020]{06A07, 05D05}
\keywords{Boolean lattice, induced poset saturation, saturated family}

\begin{document}

\begin{abstract}
Let $\Nposet$ be the four-element poset with relations
$a<b$, $c<b$, and $c<d$. We prove
\[
\satstar(n,\Nposet)=2n \qquad (1\leq n\leq 5).
\]
The cases $n=1,2$ are immediate. For $n=3,4,5$, the proof is a
complete exact enumeration after reducing induced saturation to maximal
independence in a $4$-uniform hypergraph. The numbers of labeled minimum
families are $9$, $60$, and $450$; up to ground-set permutations and
complementation, they form $2$, $4$, and $7$ classes. Thus Conjecture~16 of
Ivan and Wang is verified in the first three nonvacuous dimensions. No claim
is made for $n\geq 6$.
\end{abstract}

\maketitle

\section{Result}

Let $\Bn=2^{[n]}$, ordered by inclusion. A family
$\mathcal F\subseteq\Bn$ is \emph{induced-$\Nposet$-saturated} if it contains
no induced copy of $\Nposet$, while $\mathcal F\cup\{S\}$ contains one for
every $S\in\Bn\setminus\mathcal F$. The minimum size of such a family is
$\satstar(n,\Nposet)$.

Ferrara et al. established the upper bound
\[
\satstar(n,\Nposet)\leq 2n \qquad (n\geq 3)
\]
by an explicit construction \cite[Proposition~2]{FerraraEtAl}. Ivan and Wang
proved $\satstar(n,\Nposet)\geq(n+6)/4$ \cite[Theorem~12]{IvanWang}. They
conjectured that the upper bound is always exact
\cite[Conjecture~16]{IvanWang}.

The nearest neighbouring case is the four-element poset $2C_2$, two disjoint
$2$-chains: there the upper bound $2n$ is likewise known, and Martin and Veldt
raised the lower bound to $3n/2+1/2$ and exhibited many induced-$2C_2$-saturated
families of size $2n$ \cite{MartinVeldt}. For $\Nposet$ the present note
determines the value exactly in the smallest dimensions instead of bounding it.

\begin{theorem}\label{thm:main}
For every $1\leq n\leq 5$,
\[
\satstar(n,\Nposet)=2n.
\]
For $n=3,4,5$, the numbers of labeled minimum families are, respectively,
\[
9,\qquad 60,\qquad 450.
\]
Under ground-set permutations and complementation, these minimum families
form, respectively, $2$, $4$, and $7$ classes.
\end{theorem}

The theorem is a finite verification of the Ivan--Wang conjecture, not a proof
of its general form.

\section{Exact reduction}

For a collection $Q$ of subsets, let $G(Q)$ be its comparability graph: two
vertices are adjacent exactly when the corresponding sets are comparable by
inclusion.

\begin{lemma}\label{lem:basic}
Every induced-$\Nposet$-saturated family in $\Bn$ contains $\varnothing$ and
$[n]$. Moreover, four distinct subsets induce $\Nposet$ if and only if their
comparability graph is a path $P_4$.
\end{lemma}

\begin{proof}
Every element of $\Nposet$ is incomparable with another element. Since
$\varnothing$ and $[n]$ are comparable with every member of $\Bn$, neither can
belong to an induced copy of $\Nposet$. If either were absent from a saturated
family, adding it would create no induced copy, a contradiction.

The comparability graph of $\Nposet$ is $P_4$. Conversely, orient every edge
of a comparability graph $P_4$ from the smaller set to the larger set. Two
consecutive edges cannot point in the same direction, since transitivity would
create a chord. Hence the orientation alternates, giving $\Nposet$ or its
dual. These two posets are isomorphic.
\end{proof}

Set
\[
V_n=\Bn\setminus\{\varnothing,[n]\},
\]
and let $\mathcal H_n=(V_n,\mathcal E_n)$ be the $4$-uniform hypergraph whose
edges are the four-subsets inducing $\Nposet$.

\begin{lemma}\label{lem:hypergraph}
For $X\subseteq V_n$, the family $\{\varnothing,[n]\}\cup X$ is
induced-$\Nposet$-saturated if and only if $X$ is a maximal independent set of
$\mathcal H_n$.
\end{lemma}

\begin{proof}
By Lemma~\ref{lem:basic}, the two extreme sets cannot occur in an induced copy
of $\Nposet$. Thus $\{\varnothing,[n]\}\cup X$ is induced-$\Nposet$-free
exactly when $X$ contains no edge of $\mathcal H_n$. For $v\in V_n\setminus
X$, adding $v$ creates an induced copy exactly when some edge of $\mathcal
H_n$ is contained in $X\cup\{v\}$. This is precisely maximal independence.
\end{proof}

The number of hyperedges also has a closed form, providing an independent
check on their generation.

\begin{lemma}\label{lem:edgecount}
For every $n\geq 1$,
\[
|\mathcal E_n|=8^n-3\cdot 7^n+3\cdot 6^n-5^n.
\]
\end{lemma}

\begin{proof}
First, an induced copy of $\Nposet$ determines its roles uniquely, because
$\operatorname{Aut}(\Nposet)$ is trivial: an automorphism preserves the set of
minimal elements $\{a,c\}$ and the number of upper covers, and $a$ is the
unique minimal element with one upper cover while $c$ is the unique minimal
element with two, so $a$ and $c$ are fixed, and then $b$ (the upper cover of
$a$) and $d$ (the remaining element) are fixed as well. Hence a four-element
subset of $\Bn$ inducing $\Nposet$ admits exactly one labelling $A,B,C,D$ with
$A<B$, $C<B$, and $C<D$, and counting induced copies of $\Nposet$ is the same
as counting ordered quadruples $(A,B,C,D)$ of subsets of $[n]$ that induce
$\Nposet$ in these roles. By Lemma~\ref{lem:basic} no induced copy contains
$\varnothing$ or $[n]$, so all four sets lie in $V_n$ and this common count is
$|\mathcal E_n|$.

So let $(A,B,C,D)$ satisfy $A\subseteq B$, $C\subseteq B$, and $C\subseteq D$.
For each coordinate, the membership pattern
$(\mathbf 1_A,\mathbf 1_B,\mathbf 1_C,\mathbf 1_D)$ is then one of
\[
0000,\ 0001,\ 0100,\ 0101,\ 1100,\ 1101,\ 0111,\ 1111,
\]
these being exactly the eight of the sixteen patterns compatible with the three
inclusions. The quadruple induces $\Nposet$ in the roles $A,B,C,D$ if and only
if each of $0001$, $1100$, and $0111$ occurs.
Indeed, these patterns respectively witness $D\nsubseteq B$,
$A\nsubseteq D$, and $C\nsubseteq A$; together with the three prescribed
inclusions, they also force all strictness and all remaining
incomparabilities. Inclusion--exclusion over the three required patterns gives
the formula.
\end{proof}

\section{Enumeration and symmetry classes}

For a triple $T\in\binom{V_n}{3}$, define its completion mask
\[
M_n(T)=\{v\in V_n:T\cup\{v\}\in\mathcal E_n\}.
\]
For each candidate $X\subseteq V_n$, the verifier forms
\[
C_n(X)=\bigcup_{T\in\binom{X}{3}}M_n(T).
\]
Then $X$ is independent exactly when $C_n(X)\cap X=\varnothing$, and it is
maximal independent exactly when, in addition,
$C_n(X)\cup X=V_n$.

For each $n\in\{3,4,5\}$, every candidate with $|X|\leq 2n-2$ was tested.
The number tested was
\[
\sum_{k=0}^{2n-2}\binom{2^n-2}{k}.
\]
All operations were exact integer and bit-mask operations.

The enumeration records the size of every maximal independent set it finds.
Writing $m_n(k)$ for the number of maximal independent sets of $\mathcal H_n$
with $|X|=k$, the tested range $0\leq k\leq 2n-2$ gives
\begin{equation}\label{eq:bysize}
\begin{aligned}
\bigl(m_3(k)\bigr)_{k=0}^{4}&=(0,0,0,0,9),\\
\bigl(m_4(k)\bigr)_{k=0}^{6}&=(0,0,0,0,0,0,60),\\
\bigl(m_5(k)\bigr)_{k=0}^{8}&=(0,0,0,0,0,0,0,0,450).
\end{aligned}
\end{equation}
In particular $m_n(k)=0$ for every $k<2n-2$: no maximal independent set of
$\mathcal H_n$ is smaller than $2n-2$. The final entries $9$, $60$, $450$ are
the totals in the last column of Table~\ref{tab:census}.

\begin{table}[ht]
\centering
\begin{tabular}{@{}rrrrrr@{}}
\toprule
$n$ & $|V_n|$ & $|\mathcal E_n|$ & tested & independent & maximal \\
\midrule
$3$ & $6$  & $6$      & $57$           & $51$          & $9$   \\
$4$ & $14$ & $156$    & $6\,476$       & $3\,110$      & $60$  \\
$5$ & $30$ & $2\,550$ & $8\,656\,937$ & $1\,690\,127$ & $450$ \\
\bottomrule
\end{tabular}
\caption{Complete enumeration for $|X|\leq 2n-2$. Every maximal independent
set counted in the last column has size exactly $2n-2$; the sizes are broken
out in \eqref{eq:bysize}.}
\label{tab:census}
\end{table}

\begin{proof}[Proof of Theorem~\ref{thm:main}]
For $n=1,2$, the lattice $\Bn$ contains no induced copy of $\Nposet$, so its
only induced-$\Nposet$-saturated family is $\Bn$ itself. Hence
$\satstar(1,\Nposet)=2$ and $\satstar(2,\Nposet)=4$.

Now let $n\in\{3,4,5\}$. By Lemma~\ref{lem:hypergraph}, an
induced-$\Nposet$-saturated family of size at most $2n$ corresponds exactly to
a maximal independent set $X\subseteq V_n$ with $|X|\leq 2n-2$. By
\eqref{eq:bysize} there is none with $|X|<2n-2$, so $\satstar(n,\Nposet)\geq
2n$, and the number with $|X|=2n-2$ is $9$, $60$, and $450$ respectively.

The upper bound is exhibited on this range by the enumeration itself: each
representative $X$ of Table~\ref{tab:orbits} has $|X|=2n-2$, so
$\{\varnothing,[n]\}\cup X$ is an induced-$\Nposet$-saturated family of size
$2n$ by Lemma~\ref{lem:hypergraph}. It also follows for all $n\geq 3$ from the
construction of Ferrara et al.\ \cite[Proposition~2]{FerraraEtAl}.
\end{proof}

It remains to record the symmetry classification. Let
\[
\Gamma_n=S_n\times C_2,
\]
where $S_n$ permutes the ground set and the nontrivial element of $C_2$ sends
each subset to its complement. Complementation reverses inclusion but
preserves induced-$\Nposet$-saturation because $\Nposet$ is self-dual.

To define canonical representatives, encode a subset $S\subseteq[n]$ by
\[
\mu(S)=\sum_{i\in S}2^{i-1},
\]
sort each family by increasing $\mu$, and take the lexicographically least
image under $\Gamma_n$. In Table~\ref{tab:orbits}, a word such as \texttt{125}
denotes $\{1,2,5\}$; only the members of $X$ are displayed.

\begin{table}[ht]
\centering
\small
\begin{tabular}{@{}crl@{}}
\toprule
$n$ & orbit size & canonical representative $X$ \\
\midrule
$3$ & $6$   & \texttt{1,2,12,3} \\
$3$ & $3$   & \texttt{1,2,13,23} \\
\addlinespace
$4$ & $24$  & \texttt{1,2,12,3,123,4} \\
$4$ & $6$   & \texttt{1,2,12,3,4,34} \\
$4$ & $24$  & \texttt{1,2,12,3,124,34} \\
$4$ & $6$   & \texttt{1,2,12,123,124,34} \\
\addlinespace
$5$ & $120$ & \texttt{1,2,12,3,123,4,1234,5} \\
$5$ & $60$  & \texttt{1,2,12,3,123,4,5,45} \\
$5$ & $120$ & \texttt{1,2,12,3,123,4,1235,45} \\
$5$ & $60$  & \texttt{1,2,12,3,123,1234,1235,45} \\
$5$ & $30$  & \texttt{1,2,12,3,4,34,1234,5} \\
$5$ & $30$  & \texttt{1,2,12,3,4,34,125,345} \\
$5$ & $30$  & \texttt{1,2,12,123,1234,1235,45,345} \\
\bottomrule
\end{tabular}
\caption{Minimum families up to ground-set permutations and complementation.}
\label{tab:orbits}
\end{table}

\begin{proposition}\label{prop:orbits}
For $n=3,4,5$, the minimum families have the orbit representatives and orbit
sizes listed in Table~\ref{tab:orbits}. In particular, there are $2$, $4$, and
$7$ orbits, respectively.
\end{proposition}

\begin{proof}
For each minimum family, the classifier applies all $2\cdot n!$ elements of
$\Gamma_n$ and selects the canonical image defined above. Grouping equal
images gives Table~\ref{tab:orbits}. The orbit sizes sum to $9$, $60$, and
$450$, agreeing with the complete labeled census.
\end{proof}

This is a classification under the natural symmetries and duality of the
Boolean lattice, not under arbitrary abstract poset isomorphism. It gives a
finite partial answer to Ivan and Wang's structural Question~15
\cite{IvanWang}.

\section*{Computational specification}

Nothing outside this note is needed. The computation is specified completely
by the following five steps, each a finite exact computation that a reader can
reimplement from this description alone; Tables~\ref{tab:census}
and~\ref{tab:orbits} and \eqref{eq:bysize} are its entire output.

\begin{enumerate}
\item \emph{Ground set.} Identify $S\subseteq[n]$ with the bit-mask $\mu(S)$,
so that $V_n$ is the set of integers $1,\dots,2^n-2$ and $S\subseteq T$ holds
exactly when the bitwise conjunction of $\mu(S)$ and $\mu(T)$ is $\mu(S)$.

\item \emph{Edges.} For every $4$-subset of $V_n$, decide membership in
$\mathcal E_n$ by testing whether its comparability graph is a path on four
vertices (Lemma~\ref{lem:basic}); equivalently, by testing the $24$ bijections
onto $\Nposet$ for an isomorphism of induced subposets. The resulting count is
compared with Lemma~\ref{lem:edgecount}.

\item \emph{Masks.} For each $T\in\binom{V_n}{3}$, store $M_n(T)$ as a bit-mask
over $V_n$; since $|V_n|\leq 30$, one machine word suffices.

\item \emph{Census.} Enumerate every $X\subseteq V_n$ with $|X|\leq 2n-2$,
accumulate $C_n(X)$ as the disjunction of the masks of the triples of $X$, and
record $X$ as independent when $C_n(X)\cap X=\varnothing$ and as maximal
independent when in addition $C_n(X)\cup X=V_n$. This produces the counts of
Table~\ref{tab:census} and the size distribution \eqref{eq:bysize}.

\item \emph{Classes.} Realise $\Gamma_n$ as $2\cdot n!$ permutations of $V_n$: a
ground-set permutation permutes the bits of $\mu$, and complementation is
$\mu\mapsto 2^n-1-\mu$. Apply all of them to each recorded minimum family, sort
each image by increasing $\mu$, and keep the lexicographically least word.
Families with equal canonical images form the classes of
Table~\ref{tab:orbits}.
\end{enumerate}

Everything is exact integer and bit-mask arithmetic. There is no randomization,
floating-point arithmetic, external solver, or tolerance parameter, and no step
depends on the order of enumeration.

\section*{Exact verification}

Every value printed above was recomputed independently: a second program,
written from the definitions of this note rather than from the first
implementation, was run on a separate machine. It regenerates the induced
copies of $\Nposet$ in $\Bn$ directly from the definition, confirms the closed
form and the role uniqueness of Lemma~\ref{lem:edgecount}, reproduces the
census of Table~\ref{tab:census} entire by two independent methods (an
unpruned scan and a pruned depth-first search), with all $57$, $6\,476$ and
$8\,656\,937$ candidates tested and nothing sampled or truncated, checks each
of the $9$, $60$ and $450$ minimum families against the definition of
induced-$\Nposet$-saturation, and recovers the orbits and orbit sizes of
Table~\ref{tab:orbits} by closing the census under $\Gamma_n$ instead of by
canonical forms. All $105$ of its checks passed and every recomputed value
agreed with the value printed here. For $n\leq 4$ it also obtains
$\satstar(n,\Nposet)$ with no lemma at all, by scanning all $2^{2^n}$
subfamilies of $\Bn$; that scan is out of reach at $n=5$, where the passage
from the census to $\satstar(5,\Nposet)=10$ therefore rests on
Lemma~\ref{lem:hypergraph}, whose computational content, that no induced copy
of $\Nposet$ meets $\varnothing$ or $[n]$, is checked at $n=5$ and whose full
set equality is checked at $n=3,4$, and nothing is checked for $n\geq 6$
beyond the edge count of Lemma~\ref{lem:edgecount}. Both programs and the
transcript of that run are retained in an auxiliary archive and are available
on request; the specification above is what a reader should reimplement. An
independent recomputation that agrees is evidence; the proof is the enumeration
specified above.

\section*{Status and scope}

Theorem~\ref{thm:main} settles Conjecture~16 of Ivan and Wang for
$n\leq 5$ and leaves it open for every $n\geq 6$; it gives no information about
the asymptotic behaviour of $\satstar(n,\Nposet)$, about the lower bound of
\cite[Theorem~12]{IvanWang}, or about posets other than $\Nposet$.
Proposition~\ref{prop:orbits} classifies under $\Gamma_n$ only. We have not
found a previous determination of $\satstar(n,\Nposet)$ for $3\leq n\leq 5$:
a search of the poset-saturation literature, including
\cite{FerraraEtAl,IvanWang,MartinVeldt}, turned up the bounds recalled at the
start of this note but no exact value in these dimensions. That search was not
exhaustive and has not been through external review, so any earlier
determination we have missed takes precedence over this note.

\begin{thebibliography}{9}

\bibitem{FerraraEtAl}
M.~Ferrara, B.~Kay, L.~Kramer, R.~R. Martin, B.~Reiniger, H.~C. Smith,
and E.~Sullivan,
\emph{The saturation number of induced subposets of the Boolean lattice},
Discrete Math. \textbf{340} (2017), no.~10, 2479--2487,
\href{https://doi.org/10.1016/j.disc.2017.06.010}
{doi:10.1016/j.disc.2017.06.010}.

\bibitem{IvanWang}
M.-R.~Ivan and N.~Wang,
\emph{Linear saturation for $\mathcal N$ via butterflies},
arXiv:2511.08965v2 [math.CO], 2026.

\bibitem{MartinVeldt}
R.~R. Martin and N.~Veldt,
\emph{Induced saturation of the poset $2C_2$},
Ann. Comb. \textbf{30} (2026), no.~1, 77--90,
\href{https://doi.org/10.1007/s00026-025-00764-z}
{doi:10.1007/s00026-025-00764-z}; arXiv:2408.14648.

\end{thebibliography}

\end{document}